\chapter{Gật Đầu Cho Có}
\chapterintrobi{Hiểu chưa thì chưa rõ, nhưng phản xạ gật rất chuyên nghiệp}{Nod without understanding}{Nhiều bạn coi cái gật đầu như một chiếc áo khoác tạm cho sự ngại hỏi lại.}

\begin{lucbatbox}{Lục bát: Dạ vâng rất mượt}
Thầy vừa hỏi hiểu chưa nào\
Đầu em gật xuống thanh tao lạ thường\
Trong kia kiến thức còn sương\
Ngoài này biểu cảm đã dường rất thông\
Tới khi làm mới thấy không\
Cái đầu gật hộ mà lòng chưa hay
\end{lucbatbox}

\begin{tutuyetbox}{Thất ngôn tứ tuyệt: Dạ rồi để đó}
Cô hỏi một câu, em gật liền\
Nụ cười phối hợp cũng rất duyên\
Chỉ tiếc phần hiểu đi hơi chậm nhịp\
Nên bài về nhà vẫn đảo điên
\end{tutuyetbox}

\begin{ghichu}
Gật đầu cho có giúp tiết học trôi êm, nhưng thường đẩy phần khó sang buổi tối với mức độ khó chịu lớn hơn.
\end{ghichu}

\begin{englishnote}
\textbf{nod}: gật đầu.\par
\textbf{understand}: hiểu.\par
\textbf{I nodded, but I was still confused.}: Em gật đầu nhưng vẫn còn rối.
\end{englishnote}